\subsection{BÀI TẬP TRẮC NGHIỆM}
\ind{PHẦN I.} \inden{Câu trắc nghiệm nhiều phương án lựa chọn.  Mỗi câu hỏi học sinh chỉ chọn một phương án.}\\
\setcounter{ex}{0}
\Opensolutionfile{ans}[ans/1D3-B1-1]
\begin{ex}%[1D1Y1-1]
	Khảo sát thời gian tập thể dục của một số học sinh khối $11$ thu được mẫu số liệu ghép nhóm sau:
	\begin{center}
		\begin{tabular}{|c|c|c|c|c|c|}
			\hline Thời gian & {$[0 ; 20)$} & {$[20 ; 40)$} & {$[40 ; 60)$} & {$[60 ; 80)$} &  {$[80 ;100)$} \\
			\hline Số học sinh & $5$ & $9$ & $12$ & $10$ & $6$  \\
			\hline
		\end{tabular}
	\end{center}
	Mẫu số liệu ghép nhóm đã cho có tất cả bao nhiêu nhóm?
	\choice
	{$3$}
	{$4$}
	{\True $5$}
	{$6$}
	\loigiai{
		Dựa vào bảng số liệu, mẫu số liệu ghép nhóm đã cho có tất cả $5$ nhóm.}
\end{ex}

\begin{ex}%[1D1Y1-1]
	Khảo sát thời gian tập thể dục của một số học sinh khối $11$ thu được mẫu số liệu ghép nhóm sau:		
	\begin{center}
		\begin{tabular}{|c|c|c|c|c|c|}
			\hline Thời gian & {$[0 ; 20)$} & {$[20 ; 40)$} & {$[40 ; 60)$} & {$[60 ; 80)$} &  {$[80 ;100)$} \\
			\hline Số học sinh & $5$ & $9$ & $12$ & $10$ & $6$  \\
			\hline
		\end{tabular}
	\end{center}
	Cỡ mẫu của mẫu số liệu ghép nhóm trên là bao nhiêu?
	\choice
	{$5$}
	{$12$}
	{$20$}
	{\True $42$}
	\loigiai{
		Cỡ mẫu của mẫu số liệu ghép nhóm trên là $n=5+9+12+10+6=42$.}
\end{ex}

\begin{ex}%[1D1Y1-1]
	Khảo sát thời gian tập thể dục của một số học sinh khối $11$ thu được mẫu số liệu ghép nhóm sau
	\begin{center}
		\begin{tabular}{|c|c|c|c|c|c|}
			\hline Thời gian & {$[0 ; 20)$} & {$[20 ; 40)$} & {$[40 ; 60)$} & {$[60 ; 80)$} &  {$[80 ;100)$} \\
			\hline Số học sinh & $5$ & $9$ & $12$ & $10$ & $6$  \\
			\hline
		\end{tabular}
	\end{center}
	Độ dài của mỗi nhóm trong mẫu số liệu ghép nhóm trên là bao nhiêu?
	\choice
	{$5$}
	{$12$}
	{\True $20$}
	{$42$}
	\loigiai{
		Độ dài của mỗi nhóm là hiệu số giữa đầu mút phải và đầu mút trái của mỗi nhóm.\\
		Do đó độ dài của mỗi nhóm trong mẫu số liệu trên là $20$.}
\end{ex}


\begin{ex}%[1D1Y1-1]
	Khảo sát thời gian tập thể dục của một số học sinh khối $11$ thu được mẫu số liệu ghép nhóm sau:	
	\begin{center}
		\begin{tabular}{|c|c|c|c|c|c|}
			\hline Thời gian & {$[0 ; 20)$} & {$[20 ; 40)$} & {$[40 ; 60)$} & {$[60 ; 80)$} &  {$[80 ;100)$} \\
			\hline Số học sinh & $5$ & $9$ & $12$ & $10$ & $6$  \\
			\hline
		\end{tabular}
	\end{center}
	Giá trị đại diện của nhóm $\left[40\,;\,60\right)$ là
	\choice
	{$0$}
	{$42$}
	{\True $50$}
	{$100$}
	\loigiai{
		Giá trị đại diện của nhóm $\left[40\,;\,60\right)$ là $\dfrac{40+60}{2}=50$.}
\end{ex}

\begin{ex}%[1D1Y2-1]
	Mức thưởng tết (triệu đồng) cho các nhân viên của một công ty được thống kê trong bảng sau
	\begin{center}
		\begin{tabular}{|c|c|c|c|c|c|}
			\hline Mức thưởng tết & {$[5 ; 10)$} & {$[10 ; 15)$} & {$[15 ; 20)$} & {$[20 ; 25)$} &  {$[25 ;30)$} \\
			\hline Số nhân viên & $13$ & $35$ & $47$ & $25$ & $10$  \\
			\hline
		\end{tabular}
	\end{center}
	Giá trị đại diện của nhóm $\left[15\,;\,20\right)$ là
	\choice
	{$5$}
	{\True $17,5$}
	{$30$}
	{$130$}
	\loigiai{
		Giá trị đại diện của nhóm $\left[15\,;\,20\right)$ là $\dfrac{15+20}{2}=17,5$.}
\end{ex}

\begin{ex}%[1D1Y1-1]
	Khi thống kê chiều cao của học sinh khối lớp $12$ trong một trường trung học, ta thu được mẫu số liệu ghép nhóm sau
	\begin{center}
		\begin{tabular}{|c|c|c|c|c|c|c|}
			\hline Chiều cao $(\mathrm{cm})$ & {$[150 ; 156)$} & {$[156 ; 162)$} & {$[162 ; 168)$} & {$[168 ; 174)$} & {$[174 ; 180)$} & {$[180 ; 186)$} \\
			\hline Số học $\sinh$ & $5$ & $18$ & $40$ & $26$ & $8$ & $3$ \\
			\hline
		\end{tabular}
	\end{center}
	Độ dài của mỗi nhóm trong mẫu số liệu ghép nhóm trên là bao nhiêu?
	\choice
	{$3$}
	{\True $6$}
	{$40$}
	{$100$}
	\loigiai{
		Độ dài của mỗi nhóm là hiệu số giữa đầu mút phải và đầu mút trái của mỗi nhóm.\\
		Do đó độ dài của mỗi nhóm trong mẫu số liệu trên là $6$.}
\end{ex}

\begin{ex}%[1D1Y1-1]
	Khi thống kê chiều cao của học sinh khối $12$ trong một trường trung học, ta thu được mẫu số liệu ghép nhóm sau
	\begin{center}
		\begin{tabular}{|c|c|c|c|c|c|c|}
			\hline Chiều cao $(\mathrm{cm})$ & {$[150 ; 156)$} & {$[156 ; 162)$} & {$[162 ; 168)$} & {$[168 ; 174)$} & {$[174 ; 180)$} & {$[180 ; 186)$} \\
			\hline Số học $\sinh$ & $5$ & $18$ & $40$ & $26$ & $8$ & $3$ \\
			\hline
		\end{tabular}
	\end{center}
	Có bao học sinh khối $12$ có chiều cao từ $162$ cm đến dưới $168$ cm?
	\choice
	{$3$}
	{$6$}
	{\True $40$}
	{$100$}
	\loigiai{
		Số học sinh khối $12$ có chiều cao từ $162$ cm đến dưới $168$ cm là $40$ học sinh.}
\end{ex}


\begin{ex}%[1D1B1-4]
	Thời gian truy cập Internet mỗi buổi tối của một số học sinh được cho trong bảng sau:
	\begin{center}
		\begin{tabular}{|c|c|c|c|c|c|}
			\hline Thời gian (phút) & {$[9{,}5 ; 12{,}5)$} & {$[12{,}5 ; 15{,}5)$} & {$[15{,}5 ; 18{,}5)$} & {$[18{,}5 ; 21{,}5)$} & {$[21{,}5 ; 24{,}5)$} \\
			\hline Số học sinh & $3$ & $12$ & $15$ & $24$ & $2$ \\
			\hline
		\end{tabular}
	\end{center}
	Có bao nhiêu học sinh truy cập Internet mỗi buổi tối có thời gian từ $18{,}5$ phút đến dưới $21{,}5$ phút?
	\choice{\True $24$}
	{$15$}
	{$2$}
	{$20$}
	\loigiai{
		Số học sinh truy cập Internet mỗi buồi tối có thời gian từ $18{,}5$ phút đến dưới $21{,}5$ phút là $24$.
	}
\end{ex}

\begin{ex}%[1D2T1-2]
	Công ty $X$ trúng thầu may đồng phục cho học sinh trường $Y$ và họ đã đưa ra cỡ quần tây dựa vào số đo vòng eo của học sinh nam như sau:
	$$	\begin{tabular}{|c|c|c|c|c|c|c|}
		\hline
		Vòng eo	& $[60;64)$ &  $[64;68)$ &  $[68;72)$ &  $[72;76)$ &  $[76;80)$ \\
		\hline
		Cỡ quần	& $S$ & $M$ & $L$ & $XL$ & $XXL$ \\
		\hline
	\end{tabular}\\$$
	Và nhà trường đã gửi số đo của 40 học sinh của một lớp để công ty ước lượng tỉ lệ các cỡ quần khi may cho học sinh lớp 11
	$$\begin{tabular}{|c|c|c|c|c|c|c|c|c|c|}
		\hline
		64	& 67 & 74 & 71 & 60 & 64 & 73 & 79 & 66 & 69 \\
		\hline
		63	& 65 & 67 & 75 & 78 & 77 & 71 & 71 & 63 & 60 \\
		\hline
		61	& 65 & 71 & 71 & 66 & 73 & 63 & 77 & 65 & 63 \\
		\hline
		74	& 71 & 62 & 67 & 67 & 72 & 71 & 71 & 79 & 66 \\
		\hline
	\end{tabular}$$
	Công ty may 700 quần áo đồng phục cho học sinh lớp 11 thì nên may số lượng quần áo cỡ $S$ là bao nhiêu chiếc?
	\choice
	{129}
	{137}
	{\True140}
	{153}
	\loigiai{
		Ta có bảng số liệu ghép lớp sau:
		$$	\begin{tabular}{|c|c|c|c|c|c|c|}
			\hline
			Vòng eo	& $[60;64)$ &  $[64;68)$ &  $[68;72)$ &  $[72;76)$ &  $[76;80)$ \\
			\hline
			Cỡ quần	& $S$ & $M$ & $L$ & $XL$ & $XXL$ \\
			\hline
			Số học sinh	& 8 & 12 & 9 & 6 & 5 \\
			\hline
		\end{tabular}\\$$
		Tỉ lệ học sinh mặc vừa cỡ $S$ là: $(8:40).100\%=20\%$ nên số lượng quần cỡ $S$ công ty nên may là: $700.20\%=140$.}
\end{ex}

\begin{ex}%[1D1Y1-3]
	Khảo sát thời gian tập thể dục của một số học sinh khối $11$ thu được mẫu số liệu ghép nhóm sau:
	\begin{center}
		\begin{tabular}{|c|c|c|c|c|c|}
			\hline Thời gian & {$[0 ; 20)$} & {$[20 ; 40)$} & {$[40 ; 60)$} & {$[60 ; 80)$} &  {$[80 ;100)$} \\
			\hline Số học sinh & $5$ & $9$ & $12$ & $10$ & $6$  \\
			\hline
		\end{tabular}
	\end{center}
	Nhóm chứa mốt của mẫu số liệu ghép nhóm trên là
	\choice
	{$\left[20\,;\,40\right)$}
	{\True $\left[40\,;\,60\right)$}
	{$\left[60\,;\,80\right)$}
	{$\left[80\,;\,100\right)$}
	\loigiai{
		Nhóm chứa mốt của mẫu số liệu ghép nhóm là nhóm có tần số lớn nhất.\\
		Do đó nhóm chứa mốt của mẫu số liệu ghép nhóm trên là nhóm $\left[40\,;\,60\right)$.}
\end{ex}

\begin{ex}%[1D1Y1-3]
	Khi thống kê chiều cao của học sinh khối $12$ trong một trường trung học, ta thu được mẫu số liệu ghép nhóm sau
	\begin{center}
		\begin{tabular}{|c|c|c|c|c|c|c|}
			\hline Chiều cao $(\mathrm{cm})$ & {$[150 ; 156)$} & {$[156 ; 162)$} & {$[162 ; 168)$} & {$[168 ; 174)$} & {$[174 ; 180)$} & {$[180 ; 186)$} \\
			\hline Số học $\sinh$ & $5$ & $18$ & $40$ & $26$ & $8$ & $3$ \\
			\hline
		\end{tabular}
	\end{center}
	Nhóm chứa mốt của mẫu số liệu ghép nhóm trên là
	\choice
	{$\left[150\,;\,156\right)$}
	{\True $\left[162\,;\,168\right)$}
	{$\left[168\,;\,174\right)$}
	{$\left[180\,;\,186\right)$}
	\loigiai{
		Nhóm chứa mốt của mẫu số liệu ghép nhóm là nhóm có tần số lớn nhất.\\
		Do đó nhóm chứa mốt của mẫu số liệu ghép nhóm trên là nhóm $\left[162\,;\,168\right)$.}
\end{ex}

\begin{ex}%%[1D1Y1-2]
	Cho mẫu số liệu với cỡ mẫu $n$ được cho dưới bảng tần số ghép nhóm
	\begin{center}
		\begin{tabular}{|c|c|c|c|c|}
			\hline Nhóm & {$\left[u_1 ; u_2\right)$} & {$\left[u_2 ; u_3\right)$} & $\ldots$ & {$\left[u_k ; u_{k+1}\right)$} \\
			\hline Giá trị đại diện & $c_1$ & $c_2$ & $\ldots$ & $c_k$ \\
			\hline Tần số & $n_1$ & $n_2$ & $\ldots$ & $n_k$ \\
			\hline
		\end{tabular}
	\end{center}
	Số trung bình $\overline x $ của mẫu số liệu trên được tính bằng công thức nào sau đây
	\choice
	{$\overline x=\dfrac{u_1+u_2+\ldots+u_k}{n}$}
	{$\overline x=\dfrac{c_1+c_2+\ldots+c_k}{n}$}
	{$\overline x=\dfrac{n_1u_1+n_2u_2+\ldots+n_k{u_k}}{n}$}
	{\True $\overline x=\dfrac{n_1c_1+n_2c_2+\ldots+n_k{c_k}}{n}$}
	\loigiai{}
\end{ex}

\begin{ex}%[1D1B1-2]
	Khảo sát về cân nặng của các học sinh lớp $11D3$ người ta được một mẫu dữ liệu ghép nhóm như sau:
	\begin{center}
		\begin{tabular}{|c|c|c|c|c|c|c|}
			\hline Cân nặng & {$[30 ; 40)$} & {$[40 ; 50)$} & {$[50 ; 60)$} & {$[60 ; 70)$} & {$[70 ; 80)$} & {$[80 ; 90)$} \\
			\hline Số học sinh & $2$ & $10$ & $16$ & $8$ & $2$ & $2$ \\
			\hline
		\end{tabular}
	\end{center}
	Số trung bình của mẫu số liệu trên là
	\choice
	{\True $56$}
	{$50$}
	{$60$}
	{$55$}
	\loigiai{
		Ta có: Số phần tử của mẫu là $n=40$ và
		\begin{center}
			\begin{tabular}{|c|c|c|c|c|c|c|}
				\hline Cân nặng & {$[30 ; 40)$} & {$[40 ; 50)$} & {$[50 ; 60)$} & {$[60 ; 70)$} & {$[70 ; 80)$} & {$[80 ; 90)$} \\
				\hline Giá trị đại diện & $35$ & $45$ & $55$ & $65$ & $75$ & $85$\\
				\hline Số học sinh & $2$ & $10$ & $16$ & $8$ & $2$ & $2$ \\
				\hline
			\end{tabular}
		\end{center}
		Do đó giá trị trung bình của mẫu số liệu trên là\\
		$\overline x=\dfrac{35\cdot2+45\cdot10+55\cdot16+65\cdot8+75\cdot2+85\cdot2}{40}=56$.}
\end{ex}

\begin{ex}%[1D1B1-2]
	Thống kê về thời lượng mỗi trận đấu bi-a trong vòng tứ kết giải đấu European Open người ta được mẫu số liệu ghép nhóm như sau
	\begin{center}
		\begin{tabular}{|c|c|c|c|c|c|}
			\hline Thời gian & {$[9{,}5 ; 12{,}5)$} & {$[12{,}5 ; 15{,}5)$} & {$[15{,}5 ; 18{,}5)$} & {$[18{,}5 ; 21{,}5)$} & {$[21{,}5 ; 24{,}5)$} \\
			\hline Số trận & $3$ & $12$ & $15$ & $24$ & $2$ \\
			\hline
		\end{tabular}
	\end{center}
	Số trung bình của mẫu số liệu trên gần nhất với giá trị nào sau đây
	\choice{$17$}
	{\True $17{,}5$}
	{$18$}
	{$18{,}5$}
	\loigiai{
		Ta có số phần tử của mẫu là $n=56$ và
		\begin{center}
			\begin{tabular}{|c|c|c|c|c|c|}
				\hline Thời gian & {$[9{,}5 ; 12{,}5)$} & {$[12{,}5 ; 15{,}5)$} & {$[15{,}5 ; 18{,}5)$} & {$[18{,}5 ; 21{,}5)$} & {$[21{,}5 ; 24{,}5)$} \\
				\hline Giá trị đại diện & $11$ & $14$ & $17$ & $20$ & $2$3 \\
				\hline Số trận & $3$ & $12$ & $15$& $24$ & $2$ \\
				\hline
			\end{tabular}
		\end{center}
		Do đó giá trị trung bình của mẫu số liệu trên là
		$$
		\overline{x}=\dfrac{11\cdot3+14\cdot12+17\cdot15+20\cdot24+23\cdot2}{56}=\dfrac{491}{28} \approx 17{,}54.$$
	}
\end{ex}

\begin{ex}%[1D1B1-2]
	Doanh thu bán hàng trong $20$ ngày được lựa chọn ngẫu nhiên của một của hàng được ghi lại ở bảng sau (đơn vị: triệu đồng)
	\begin{center}
		\begin{tabular}{|c|c|c|c|c|c|}
			\hline Doanh thu & {$[5 ; 7)$} & {$[7 ; 9)$} & {$[9 ; 11)$} & {$[11 ; 13)$} & {$[13 ; 15)$} \\
			\hline Số ngày & $2$ & $7$ & $7$ & $3$ & $1$ \\
			\hline
		\end{tabular}
	\end{center}
	Số trung bình của mẫu số liệu trên thuộc khoảng nào trong các khoảng dưới đây?
	\choice{ $[7 ; 9)$}
	{ \True $[9 ; 11)$}
	{ $[11 ; 13)$}
	{ $[13 ; 15)$}
	\loigiai{
		Bảng tần số ghép nhóm theo giá trị đại diện là
		\begin{center}
			\begin{tabular}{|c|c|c|c|c|c|}
				\hline Doanh thu & {$[5 ; 7)$} & {$[7 ; 9)$} & {$[9 ; 11)$} & {$[11 ; 13)$} & {$[13 ; 15)$} \\
				\hline Giá trị đại diện & $6$ & $8$ & $10$ & $12$ & $14$ \\
				\hline Số ngày & $2$ & $7$ & $7$ & $3$ & $1$ \\
				\hline
			\end{tabular}
		\end{center}
		Số trung bình $\overline{x}=\dfrac{2\cdot 6+7\cdot8+7\cdot10+3\cdot12+1\cdot 14}{20}=9{,}4$.
	}
\end{ex}

\begin{ex}%[1D1B1-2]
	Trung tâm ngoại ngữ thống kê bảng điểm môn Tiếng Anh của một khóa học trong bảng bên dưới
	\begin{center}
		\begin{tabular}{|l|c|c|c|c|c|}
			\hline
			Điểm     & [0;2) & [2;4) & [4;6) & [6;8) & [8;10) \\ \hline
			Học viên & 10   & 30   & 55   & 42   & 9     \\ \hline
		\end{tabular}
	\end{center}
	Số trung bình của mẫu số liệu thuộc khoảng nào trong các khoảng dưới đây?
	\choice
	{$[8;10)$}
	{\True $[4;6)$}
	{$[2;4)$}
	{$[6;8)$}
	\loigiai{
		Ta có bảng thống kê theo giá trị đại diện như sau:
		\begin{center}
			\begin{tabular}{|l|c|c|c|c|c|}
				\hline
				Giá trị đại diện & 1  & 3  & 5  & 7  & 9 \\ \hline
				Tần số           & 10 & 30 & 55 & 42 & 9 \\ \hline
			\end{tabular}
		\end{center}
		Khi đó ta có số trung bình của mẫu số liệu trên được tính như sau: \\
		$$\bar{x}=\dfrac{1.10+3.30+5.55+7.42+9.9}{10+30+55+42+9}\approx 5,14.$$
	}
\end{ex}

\begin{ex}%[1D1B1-3]
	Khảo sát về số giờ mượn sách thư viện của học sinh khối $11$ trường Y ta được một mẫu số liệu ghép nhóm như sau
	\begin{center}
		\begin{tabular}{|c|c|c|c|c|c|}
			\hline Số giờ mượn & {$[0 ; 4)$} & {$[4 ; 8)$} & {$[8 ; 12)$} & {$[12 ; 16)$} & {$[16 ; 20)$} \\
			\hline Số học sinh & $54$ & $78$ & $120$ & $45$ & $12$\\
			\hline
		\end{tabular}
	\end{center}
	Mốt của mẫu số liệu ghép nhóm trên gần nhất với giá trị nào sau đây
	\choice
	{$120$}
	{$12$}
	{$8$}
	{\True $9$}
	\loigiai{
		Nhóm chứa Mốt là $\left[8;12\right)$. Do đó:\\
		$u_m=8$ ; $n_{m-1}=78$ ; $n_m=120$ ; $n_{m+1}=45$ ; $u_{m+1}-u_m=12-8=4$.\\
		Mốt của mẫu số liệu ghép nhóm là: $M_0=8+\dfrac{120-78}{\left(120-78\right)+\left(120-45\right)}.4=\dfrac{368}{39}\approx 9{,}4$.}
\end{ex}

\begin{ex}%[1D1B1-3]
	Khảo sát thời gian tập thể dục của một số học sinh khối $11$ thu được mẫu số liệu ghép nhóm sau
	\begin{center}
		\begin{tabular}{|c|c|c|c|c|c|}
			\hline Thời gian (phút) & {$[0 ; 20)$} & {$[20 ; 40)$} & {$[40 ; 60)$} & {$[60 ; 80)$} & {$[80 ; 100)$} \\
			\hline Số học sinh & $5$ & $9$ & $12$ & $10$ & $6$ \\
			\hline
		\end{tabular}
	\end{center}
	Mốt của mẫu số liệu trên là
	\choice{$42 $}
	{\True  $52$}
	{ $53$}
	{ $54$}
	\loigiai{
		Mốt $M_0$ chứa trong nhóm $[40 ; 60)$.\\
		Do đó $u_m=40$, $u_{m+1}=60 \Rightarrow u_{m+1}-u_m=60-40=20$.\\
		$n_{m-1}=9$, $n_m=12$, $n_{m+1}=10$.
		$$ M_0=40+\dfrac{12-9}{(12-9)+(12-10)}(60-40)=52.$$
	}
\end{ex}

\Closesolutionfile{ans}

\ind{PHẦN II.} \inden{Câu trắc nghiệm đúng sai. Trong mỗi ý a), b), c), d) ở mỗi câu, học sinh chọn đúng hoặc sai.}\\
\Opensolutionfile{ans}[ans/1D3-B1-2]


\begin{ex}
	Cho mẫu số liệu ghép nhóm về điểm thi và số người dự thi như sau:
	
	\begin{center}\begin{tabular}{|c|c|c|c|c|c|c|}
			\hline
			Điểm thi   & [2 ; 4) & [4 ; 6) & [6 ; 8) & [8 ; 10) & [10 ; 12)\\  
			\hline 
			Số người dự thi & 4 & 20 & 9 & 16 & 15 \\ 
			\hline 
		\end{tabular}
	\end{center}
	\choiceTF
	{\True Độ dài mỗi nhóm của mẫu số liệu ghép nhóm trên là $2$}
	{Cỡ mẫu của mẫu số liệu ghép nhóm trên là $62$}
	{\True Số trung bình của mẫu số liệu ghép nhóm trên là ${7,56}$ }
	{\True Mốt của mẫu số liệu ghép nhóm trên là $\dfrac{140}{27}$}

	\loigiai{ 
		\begin{enumerate}[a)]
			\item Độ dài mỗi nhóm là $4-2=2$
			\item Cõ mẫu $4+20+9+16+15=64$
			\item Các giá trị đại diện của mẫu số liệu là: 3; 5; 7; 9; 11
			
			Tổng tần số là: $n=64$
			
			Số trung bình của mẫu số liệu ghép nhóm là:
			
			$\overline{x}=\dfrac{3.4+5.20+7.9+9.16+11.15}{ 64}=7,56$.
			\item Nhóm chứa mốt là $[4;6)$. Suy ra
			$$M_o=4+\dfrac{20-4}{(20-4)+(20-9)}.2=\dfrac{140}{27}$$
		\end{enumerate}

		
}\end{ex}

% \begin{ex}%[1D3H1-3]
% 	Kết quả khảo sát cân nặng của $25$ quả cam ở mỗi lô hàng $A,$ $ B$ được cho ở bảng sau:
% 	\begin{center}
% 		\begin{tabular}{|c|c|c|c|c|c|}
% 			\hline Cân nặng (gam) & {$[150 ; 155)$} & {$[155 ; 160)$} & {$[160 ; 165)$} & {$[165 ; 170)$} & {$[170 ; 175)$} \\
% 			\hline Số quả cam ở lô hàng $A$ & $2$ & $6$ & $12$ & $4$ & $1$ \\
% 			\hline Số quả cam ở lô hàng $B$ & $1$ & $3$ & $7$ & $10$ & $4$ \\
% 			\hline
% 		\end{tabular}
% 	\end{center}
% 	\choiceTF
% 	{\True Giá trị đại diện nhóm $[150;155)$ bằng $152{,}5$}
% 	{ Cân nặng trung bình của mỗi quả cam ở lô $A$ là $163{,}7$ (gam)}
% 	{ Cân nặng trung bình của mỗi quả cam ở lô $B$ là $162{,}1$ (gam)}
% 	{\True  Theo số trung bình thì cam ở lô hàng $B$ nặng hơn cam ở lô hàng $A$}
% 	\loigiai{
% 		\begin{itemchoice}
% 			\itemch Giá trị đại diện nhóm $[150;155)$ bằng $\dfrac{150+155}{2} =152{,}5.$
% 			\itemch Bảng thống kê số lượng cam theo giá trị đại diện:
% 			\begin{center}
% 				\begin{tabular}{|c|c|c|c|c|c|}
% 					\hline Cân nặng đại diện (gam) & $152{,}5$ & $157{,}5$ & $162{,}5$ & $167{,}5$ & $172,5$ \\
% 					\hline Số quả cam ở lô hàng $A$ & $2$ & $6$ & $12$ & $4$ & $1$ \\
% 					\hline Số quả cam ở lô hàng $B$ & $1$ & $3$ & $7$ & $10$ & $4$ \\
% 					\hline
% 				\end{tabular}
% 			\end{center}
% 			Cân nặng trung bình của mỗi quả cam ở lô $A$ là
% 			$$\bar{x}_A=\dfrac{152{,}5 \cdot 2+157{,}5 \cdot 6+162{,}5 \cdot 12+167{,}5 \cdot 4+172{,}5 \cdot 1}{25}=161{,}7 \text{ (gam).}$$
% 			\itemch Cân nặng trung bình của mỗi quả cam ở lô $B$ là
% 			$$\bar{x}_B=\dfrac{152{,}5 \cdot 1+157{,}5 \cdot 3+162{,}5 \cdot 7+167{,}5 \cdot 10+172{,}5 \cdot 4}{25}=165{,}1 \text{ (gam).}$$
% 			\itemch Ta thấy $\bar{x}_{A}<\bar{x}_{B}$. Vậy nếu so sánh theo số trung bình thì cam ở lô hàng $B$ nặng hơn cam ở lô hàng $A$.
% 	\end{itemchoice}}
% \end{ex}

\begin{ex}
	Số lượng người đi xem một bộ phim mới theo độ tuổi trong một rạp chiếu phim (sau 1 giờ đầu công chiếu) được ghi lại theo bảng phân phối ghép nhóm sau:
	\begin{center}
		\begin{tabular}{|c|c|c|c|c|c|}
			\hline Độ tuồi & {$[10 ; 20)$} & {$[20 ; 30)$} & {$[30 ; 40)$} & {$[40 ; 50)$} & {$[50 ; 60)$} \\
			\hline Số người & 6 & 12 & 16 & 7 & 2 \\
			\hline
		\end{tabular}
	\end{center}
	\choiceTF
	{\True Giá trị đại diện nhóm $[50 ; 60)$ là $55$}
	{\True Độ tuổi được dự báo là ít xem phim đó nhất là thuộc nhóm $[50;60)$}
	{\True  Nhóm chứa mốt là nửa khoảng $[30 ; 40)$}
	{Độ tuổi được dự báo là thích xem phim đó nhiều nhất là 31 tuổi}
	\loigiai{
		
	}
\end{ex}

\Closesolutionfile{ans}

\ind{PHẦN III.} \inden{Câu trắc nghiệm trả lời ngắn.}\\
\Opensolutionfile{ans}[ans/1D3-B1-3]

\begin{ex}
	Cho mẫu số liệu ghép nhóm về điểm thi và số người dự thi như sau:
	
	\begin{center}\begin{tabular}{|c|c|c|c|c|c|c|}
			\hline
			Điểm thi   & [0 ; 3,5) & [3,5 ; 7) & [7 ; 10,5) & [10,5 ; 14) & [14 ; 17,5) & [17,5 ; 21)\\  
			\hline 
			Số người dự thi & 15 & 8 & 16 & 2 & 17 & 18 \\ 
			\hline 
		\end{tabular}
	\end{center}
	Tính điểm thi trung bình từ mẫu số liệu ghép nhóm trên (kết quả làm tròn đến hàng phần chục).\\
	\shortans[oly]{$11,2$}
	\loigiai{ 
		
		Các giá trị đại diện của mẫu số liệu là: 1,75; 5,25; 8,75; 12,25; 15,75; 19,25
		
		Tổng tần số là: $n=76$
		
		Số trung bình của mẫu số liệu ghép nhóm là:
		
		$\overline{x}=\dfrac{1,75.15+5,25.8+8,75.16+12,25.2+15,75.17+19,25.18}{ 76}=11,2$.
		
		
}\end{ex}

\begin{ex}
	Cho mẫu số liệu ghép nhóm về điểm thi và số người dự thi như sau:
	
	\begin{center}\begin{tabular}{|c|c|c|c|c|c|c|}
			\hline
			Điểm thi   & [0 ; 3) & [3 ; 6) & [6 ; 9) & [9 ; 12) & [12 ; 15) & [15 ; 18)\\  
			\hline 
			Số người dự thi & 8 & 12 & 13 & 1 & 19 & 10 \\ 
			\hline 
		\end{tabular}
	\end{center}
	Tính điểm thi trung bình từ mẫu số liệu ghép nhóm trên.\\
	\shortans[oly]{$9,45$}
	\choice
	{ $ {7,95}$ }
	{ $ {9,00}$ }
	{ $ {6,00}$ }
	{ \True ${9,45}$ }
	\loigiai{ 
		
		Các giá trị đại diện của mẫu số liệu là: 1,5; 4,5; 7,5; 10,5; 13,5; 16,5
		
		Tổng tần số là: $n=63$
		
		Số trung bình của mẫu số liệu ghép nhóm là:
		
		$\overline{x}=\dfrac{1,5.8+4,5.12+7,5.13+10,5.1+13,5.19+16,5.10}{ 63}=9,45$.
		
		
}\end{ex}
\begin{ex}
	Cho mẫu số liệu ghép nhóm về điểm số và số học sinh như hình dưới đây
	\begin{center}\begin{tabular}{|c|c|c|c|c|c|c|}
			\hline
			Điểm số   & [2 ; 4,5) & [4,5 ; 7) & [7 ; 9,5) & [9,5 ; 12) & [12 ; 14,5) & [14,5 ; 17)\\  
			\hline 
			Số học sinh & 1 & 6 & 3 & 10 & 9 & 1 \\ 
			\hline 
		\end{tabular}
	\end{center}
Tìm mốt của mẫu số liệu ghép nhóm trên (kết quả làm tròn đến hàng phần chục).\\
\shortans[oly]{$11,7$}
	\loigiai{ 
		Nhóm chứa mốt là nhóm thứ $k=4, n_k=10,n_{k-1}=3, n_{k+1}=9, u_k=9.5, u_{k+1}=12.0$.
		
		Áp dụng công thức tìm Mốt: $M_o=u_k+\dfrac{ (n_k - n_{k-1}) }{(n_k-n_{k-1})+(n_k-n_{k+1}) }(u_{k+1}-u_k)=11,7$.
	
		
}\end{ex}

\begin{ex}%[1D1B1-3]
	Một công ty bất động sản khảo sát khách hàng xem họ có nhu cầu mua căn hộ ở mức giá nào. Kết quả khảo sát được ghi lại ở bảng sau:
	$$	\begin{tabular}{|c|c|c|c|c|c|}
		\hline
		Mức giá (triệu đồng/$m^2$)	& $[28;32)$ &  $[32;36)$ &  $[36;40)$ &  $[40;44)$ &  $[44;48)$ \\
		\hline
		Số khách hàng	& 40 & 78 & 120 & 51 & 11 \\
		\hline
	\end{tabular}\\$$
	Công ty nên xây nhà ở mức giá nào để nhiều người có nhu cầu mua nhất?\\
	\shortans[oly]{$37,5$}
	\loigiai{ Nhóm chứa mốt của mẫu số liệu trên là nhóm $[36;40)$.\\
		Do đó $u_m=36,n_{m-1}=78,n_m=120,n_{m+1}=51,u_{m+1}-u_m=40-36=4$.\\
		Mốt của mẫu số liệu ghép nhóm là:
		$$M_o=36+\dfrac{120-78}{(120-78)+(120-51)}.4\approx37,5.$$  }
\end{ex}

\Closesolutionfile{ans}



